% page 507 of the facsimile (image p-506 of the scan)
% block 165.1 x 237.7 mm ; left margin 36.9 mm
\pgc{15.240mm}{0.000mm}{
\l{\cel{54}499}
\l{}
\l{\sou{rigida}. I.Qua ne cedas se onu probas flexar lu. - II. (metaf.)}
\l{\cel{3}Qua ne igas flexar la reguli ; quan onu ne povas igar devia}
\l{\cel{3}car de la apliko di la reguli. - DEFIS.}
\l{}
\l{\sou{riglo.}\cel{1}Mikra peco fera, fixigita sur pordo, sur fenestro, e}
\l{\cel{3}qua, se pulsita aden seruro-boketo, impedas la aperto. -D}
\l{}
\l{\sou{rigoro.}\cel{1}Precizeso, strikteso, extrema. - DEFIS}
\l{}
\l{\sou{rikochar.}\cel{1}(\sou{netrans., sur}). Dicesas pri la retro-salto de stono}
\l{\cel{3}plata qua frapas la surfaco di aquo, o projektilo de pafarmo}
\l{\cel{3}qua frapas sur sulo. - DEFR.}
\l{}
\l{rimo. Konsonanco di la lasta silabo acentoza di la vorto finala}
\l{\cel{3}en du o plura versi. - DEFIRS.}
\l{}
\l{rimborsar. (trans., ad).\cel{2}Donar (ad ulu) la equivalanto di la}
\l{\cel{3}sumo de to quon lu preste spensis.- DEFIS.}
\l{}
\l{rimeno. I. Bendo de ledro, longa e streta. - II. Bendo de ledro}
\l{\cel{3}di qua un extremajo esas garnisita ye buklo, e la altra havas}
\l{\cel{3}longeso-sinse, serio de trui aden qui eniras la buklo-pinto}
\l{\cel{3}segun ke onu volas plektar plu o min strete. - III.Bendo kon-}
\l{\cel{3}tinua de ledro, de kauchuko, di qua la extremaji interligesis}
\l{\cel{3}tale ke formacesas cirklo-kurvo flexebla, uzata por transmisar}
\l{\cel{3}la movo di mashino. - DR.}
\l{}
\l{rinforcar. (trans.,\cel{1}\sou{per}).I. (\sou{muziko}).Igar (sono) plu intensa, plu}
\l{\cel{3}energioza. - II. (milit.) Igar plu multa la efektivo de la}
\l{\cel{3}militisti di garnizono. - EFIS.}
\l{}
\l{\sou{ringo}. I. Cirklokurvo de metalo o de altra materio rezistiva}
\l{\cel{3}quan onu uzas, en la kazi maxim multa, por retenar, ligar.}
\l{\cel{3}- II. To quo, sua-forme, similesas ringo. - III. (geom.)}
\l{\cel{3}Parto de plano, qua kontenesas inter du cirkli\cel{1}\sou{sama}-centra}
\l{\cel{3}- DE.}
\l{}
\l{rinocero. (zool.)\cel{2}Granda mamifero, sovaja, ek la klaso \textquotedbl{}\sou{paki}-}
\l{\cel{3}dermi\textquotedbl{}, kun un o du korni sur la nazo. - L. rhinoceros.- DEFIS}
\l{}
\l{rinoplastiko. Operaco kirurgiala qua konsistas ye rifacar la nazo}
\l{\cel{3}per la grefto animalala, per pelo-peco deprenita de la fronto}
\l{\cel{3}de la brakio, e c. -}
\l{}
\l{rinsar. (trans., per)\cel{2}Pasar en aquo\cel{2}pura, to quo esas jus l}
\l{\cel{3}vita, por igar lu absolute neta. - EF.}
\l{}
\l{ripar. (netrans.) Dicesas pri roto di biciklo o di automobi}
\l{\cel{3}qua glitas skrapante aden direciono qua esas obliqua ye l}
\l{\cel{3}plaso propra.- F.}
\l{}
\l{\sou{riskar}. (\sou{trans.}) Expozar (ulu, ulo) a chanco pri qua onu}
\l{\cel{3}certa (+certena) ke lu esos bona. - DEFIRS.}
}
